\subsection{STRICT INEQUALITIES BETWEEN INTEGERS}

PROPOSITION 2. Let \(a\) and \(b\) be two integers. Then \(a < b\) if and only if there exists an integer \(c > 0\) such that \(b = a + c\) .

If \(a < b\) , there exists a cardinal \(c \leqslant b\) (so that \(c\) is an integer ( \(\S 4\) , no. 2, Proposition 2)) such that \(b = a + c\) ( \(\S 3\) , no. 6, Proposition 13); if \(a \neq b\) , we must have \(c \neq 0\) . Conversely, if \(b = a + c\) and \(c \neq 0\) , then \(c \geqslant 1\) and hence \(a < a + 1 \leqslant a + c = b\) .

PROPOSITION 3. Let \((a_i)_{i \in I}\) and \((b)_{i \in I}\) be two finite families of integers such that \(a_i \leqslant b_i\) for each \(i \in I\) and \(a_i < b_i\) for at least one index \(i\) . Then

\[
\sum_ {i \in \mathbf {I}} a _ {i} <   \sum_ {i \in \mathbf {I}} b _ {i}.
\]

If also \(b_{i} > 0\) for each \(i \in \mathbf{I}\) , then

\[
\prod_ {i \in I} a _ {i} <   \prod_ {i \in I} b _ {i}.
\]

Let \(j\) be an index such that \(a_{j} < b_{j}\) , and let \(J = I - \{j\}\) . Then

\[
b _ {j} = a _ {j} + c _ {j}
\]

with \(c_{j} > 0\) (Proposition 2), and therefore (§ 3, no. 6, Proposition 14)

\[
\sum_ {i \in \mathbf {I}} b _ {i} = a _ {j} + c _ {j} + \sum_ {i \in \mathbf {J}} b _ {i} \geqslant c _ {j} + a _ {j} + \sum_ {i \in \mathbf {J}} a _ {i} = c _ {j} + \sum_ {i \in \mathbf {I}} a _ {i}.
\]

Since \(c_{j} > 0\) , the first part of the Proposition follows from Proposition 2. Likewise

\[
\prod_ {i \in I} b _ {i} = (a _ {j} + c _ {j}) \prod_ {i \in J} b _ {i} = a _ {j}. \prod_ {i \in J} b _ {i} + c _ {j}. \prod_ {i \in J} b _ {i} \geqslant \prod_ {i \in I} a _ {i} + c _ {j}. \prod_ {i \in J} b _ {i}.
\]

Since \(c_{j}\) and all the \(b_{i}\) are \(\neq 0\) , the product \(c_{j}\) . \(\prod_{i \in \mathbb{J}} b_{i}\) is \(\neq 0\) (§3, no. 4, Proposition 7); hence the second part of the Proposition.

COROLLARY 1. Let \(a, a',\) and \(b\) be integers such that \(a < a'\) and \(b > 0\) . Then \(a^b < a'^b\) .

We need only express \(a^b\) and \(a'^b\) as products of finite families of integers (§ 3, no. 5, Proposition 10) and apply Proposition 3, observing that the relation \(a < a'\) implies \(a' > 0\) .

COROLLARY 2. Let \(a, b\) , and \(b'\) be integers such that \(a > 1\) and \(b < b'\) ; then \(a^b < a^{b'}\) .

For there exists an integer \(c > 0\) such that \(b' = b + c\) (Proposition 2); since \(c \geqslant 1\) , we have \(a^c \geqslant a > 1\) , whence \(a^{b'} = a^b a^c > a^b\) .

COROLLARY 3. Let \(a, b, b'\) be integers (resp. integers such that \(a > 0\) ). Then \(a + b = a + b'\) (resp. \(ab = ab'\) ) if and only if \(b = b'\) .

COROLLARY 4. If \(a\) and \(b\) are integers such that \(a \leqslant b\) , then there exists a unique integer \(c\) such that \(b = a + c\) .

The existence of \(c\) follows from Proposition 13 of §3, no. 6, and its uniqueness from Corollary 3 above.

The integer \(c\) such that \(b = a + c\) (where \(a \leqslant b\) ) is called the difference of the integers \(b\) and \(a\) , and is written \(b - a\) . It is easily verified that if \(a, b, a', b'\) are integers such that \(a \leqslant b\) and \(a' \leqslant b'\) , then

\[
(b - a) + (b ^ {\prime} - a ^ {\prime}) = (b + b ^ {\prime}) - (a + a ^ {\prime}).
\]

